\section{Introduction}
\label{sec:introducao}

A calibração de sensores de baixo custo utilizados em monitoramento da qualidade do ar tem ganhado significativa relevância científica nos últimos anos, especialmente devido à expansão de redes urbanas e à necessidade de medições contínuas, robustas e financeiramente acessíveis. Entre esses dispositivos, os sensores do tipo Óxido Metálico (MOx) destacam-se pela ampla disponibilidade e baixo custo, embora apresentem limitações importantes, como deriva térmica, sensibilidade cruzada e forte dependência das condições meteorológicas. Para mitigar tais limitações, técnicas de regressão multivariada, aprendizado de máquina e modelos híbridos têm sido amplamente investigadas \cite{ols_wls_pls_ridge_lasso_2024,plsr_pcr_rr_mlr_pm10,unveiling_optimal_regression_2024,calibration_lowcost_no2_2019}.

Estudos recentes demonstram que técnicas clássicas, como Mínimos Quadrados Ordinários (OLS), Weighted Least Squares (WLS), Partial Least Squares (PLS), Ridge e Lasso, apresentam desempenho consistente na presença de multicolinearidade e alto ruído \cite{ols_wls_pls_ridge_lasso_2024}. Comparações sistemáticas entre PLSR, PCR, Ridge Regression e Múltipla Regressão Linear indicam que modelos multivariados capturam relações complexas entre variáveis meteorológicas e concentrações de poluentes \cite{plsr_pcr_rr_mlr_pm10}. Modelos mais avançados, como GLM, Random Forest, Multi-Layer Perceptron (MLP) e métodos de fatoração de fontes também têm sido empregados com resultados promissores \cite{unveiling_optimal_regression_2024}.

A literatura recente também revela uma tendência crescente no uso de modelos de aprendizado profundo para calibração, previsão e correção de sensores. O trabalho de \cite{calibration_lowcost_no2_2019} utiliza regressão múltipla e modelos não lineares para calibrar sensores de NO$_2$ em redes urbanas. Modelos recorrentes, como RNN e LSTM, têm sido aplicados com sucesso na previsão de concentrações internas de NO$_2$ \cite{forecasting_no2_rnn_lstm_2025}, enquanto arquiteturas híbridas MLP-LSTM alcançam alta precisão na predição de NO$_x$, NO$_2$ e C$_6$H$_6$ \cite{high_precision_nox_2024}. Abordagens combinadas, como SRA-MLP, SVM, KNN, otimizadores como PSO e modelos generativos (GAN e VAE), aparecem em estudos recentes como alternativas para lidar com não linearidades e incertezas \cite{sra_mlp_2021,improved_air_quality_ml,mlp_pso_2025,review_data_driven_2025,vae_mlp_lstm_2023}.

Neste contexto, o presente trabalho utiliza o conjunto de dados \textit{Air Quality UCI} (2004–2005), composto por 9.358 registros horários e contendo cinco sensores MOx (PT08.S1 a PT08.S5), variáveis meteorológicas (Temperatura, Umidade Relativa e Umidade Absoluta) e cinco variáveis \textit{Ground Truth} (GT), representando concentrações reais de poluentes atmosféricos (CO, NO$_x$, NO$_2$, C$_6$H$_6$, entre outros). Sensores MOx operam pela variação da resistência elétrica de semicondutores expostos a gases específicos, sendo fortemente impactados por fenômenos meteorológicos, o que exige técnicas robustas de calibração.

Para reduzir redundância entre variáveis e mitigar efeitos de multicolinearidade, aplicou-se a técnica de Análise de Componentes Principais (PCA). A transformação resultou em dois ou três componentes principais preservando entre 77,5\% e 92,1\% da variância total dos preditores, aumentando a robustez e interpretabilidade dos modelos \cite{Souza2025EDA}. Ademais, foi criado um Índice Composto de Poluição (CPI), normalizado via Z-score, permitindo integrar informações das variáveis GT em um único indicador robusto.

O objetivo deste trabalho é desenvolver e validar modelos lineares e não lineares para a calibração multivariada de sensores MOx, avaliando: (i) a precisão das leituras frente aos valores GT; (ii) o impacto das variáveis meteorológicas nas respostas dos sensores; (iii) a eficácia de técnicas de redução de dimensionalidade. A variável NMHC foi excluída devido ao elevado número de registros ausentes (superior a 90\%), garantindo a integridade estatística da análise\cite{Souza2025EDA}.

Ao combinar métodos estatísticos clássicos, aprendizado de máquina e técnicas modernas de modelagem, este estudo contribui para o avanço da calibração de sensores de baixo custo, indispensáveis para cidades inteligentes, redes participativas (\textit{citizen science}) e sistemas de monitoramento ambiental de baixo orçamento.
